% !PRINTED 568
\weberpage{568}

\begin{center}
§ 152.\qquad \textbf{Die Pole der Funktion} $\wp(\mu u)$.
\end{center}

Um den Grad der Funktionen $R$ und $P$ festzustellen, müssen wir die Nullstellen von $P$, also die Unendlichkeitsstellen von $\wp(\mu u)$, abzählen, die nicht zugleich Unendlichkeitsstellen von $\wp(u)$ sind. Es wird aber $\wp(\mu u):\wp(u)$ dann und nur dann unendlich, wenn
\weq{1}{u = \dfrac{h_1\omega_1 + h_2\omega_2}{\mu}}
wird, wenn $h_1$, $h_2$ ganze Zahlen sind, wenn $h_1\omega_1 + h_2\omega_2$, aber nicht $(h_1\omega_1+h_2\omega_2):\mu$ eine Periode ist. Setzen wir also
\weq{2}{\wp(\mu u) = \dfrac{R_\mu[\wp(u)]}{P_\mu[\wp(u)]},}
und bezeichnen mit $g$ die Wurzeln von $P_\mu(x)$, so ist
\weq{3}{g = \wp\left(\dfrac{h_1\omega_1+h_2\omega_2}{\mu}\right).}
Die Zahl $g$ ändert sich nicht, wenn $h_1$ und $h_2$ um Vielfache von $m=\mu\mu'$ geändert werden, weil dadurch nach (11), § 151 das Argument der $\wp$-Funktionen um eine Periode geändert wird.

Da die Funktionen $\wp(u)$ und $\wp(\mu u)$ nicht geändert werden, wenn wir die Periode $\omega_1$, $\omega_2$ einer linearen Transformation mit der Determinante 1 unterwerfen, so beschränken wir die Allgemeinheit nicht, wenn wir annehmen, $A$ sei relativ prim zu $m$ und wegen der Periodizität von $\wp(u)$ können wir daher auch setzen:
\[ g = \wp\left(\dfrac{h_1\omega_1+h_2 A\omega_2}{\mu}\right). \]
Wegen der Homogenität von $\wp(u)$ können wir nun $\omega_1=1$, $\omega_2=\omega$ setzen und erhalten nach § 151, (4):
\[ A\omega = \dfrac{-B+\sqrt{\Delta}}{2}, \]
und demnach ist
\weq{4}{h_1\omega_1 + h_2 A\omega_2 = h_1 + h_2\left(\dfrac{-B+\sqrt{\Delta}}{2}\right).}
Wenn wir daher
\[ h_1 + h_2\dfrac{-B+\sqrt{\Delta}}{2} = \nu \]
setzen, so ergibt sich aus (3):
\weq{5}{g = \wp\left(\dfrac{\nu}{\mu}\right),}
%CONT>>
